\section{Introduction}
\label{sec:intro}

The deployment-safety theorem of
\citep[Deployment Safety]{lasser2026paper10} establishes a
conditional capability-magnitude-independent bound on Goodhart
slack for capability-unbounded deployments under twelve
operational conditions and an explicit empirical-adequacy
assumption inherited from the
\citep[Microfoundation]{lasser2026micro} Concentration-Gap
conjecture (HHI surrogate adequacy). Two of these hypotheses sit awkwardly in the
deployment-claim's conditional structure:

\begin{enumerate}
\item \textbf{(C7) bounded co-evolution}, a free-floating
  deployment-class constant $\bar{M}$ asserted to be independent
  of capability magnitude $|P|$, with no derivation from primitive
  operational parameters. \citep[\S 10.2]{lasser2026paper10}
  identifies this as a hard-failure mode under outright failure of
  the assumption.
\item \textbf{The Concentration-Gap conjecture}, inherited from
  \citep{lasser2026micro} as deferred to follow-up work. The HHI
  surrogate-adequacy assumption serves as its operational surrogate
  but is itself an inherited empirical-adequacy claim rather than a
  derived consequence.
\end{enumerate}

Both gaps are closed by structural derivations established in the
present paper:

\begin{itemize}
\item \textbf{Theorem~\ref{thm:c7-corollary} (Bounded co-evolution
  as a corollary).} Under (C5.HOEFF) per-step LLR clipping, an
  explicit upper exposure-rate cap (C7.RATE), and the
  channel-projection structure of (C5.MULT), the per-channel
  coupling magnitude is bounded by a closed-form constant in
  primitive deployment parameters: $\bar{M}^{\mathrm{cum}} \leq C
  \cdot \Bclip \cdot \lammax \cdot \taumeta$ with $C \in \{1,
  \Kch - 1\}$. The previously free-floating bounded co-evolution
  assumption of \citep[Deployment Safety]{lasser2026paper10}
  graduates to a corollary; the property is now invoked at
  condition~(C7) of the deployment-safety theorem.

\item \textbf{Theorem~\ref{thm:concgap-scoped} (Concentration-Gap
  Selection Theorem, scoped).} Under three operationally
  auditable structural conditions --- representativeness (REP),
  bounded dispersion (DISP), coalition closure (COAL) --- and a
  Lipschitz embedding compatibility assumption
  (Assumption~\ref{ass:embedding-companion}), the proxy-truth
  Goodhart slack is bounded by a $\chi^2$-divergence quantity that
  controls $\HHI$ via $\chi^2(w \,\|\, \mu) = N \HHI(w) - 1$ for
  uniform $\mu$ on a (COAL)-bounded counterparty population. The
  HHI surrogate-adequacy assumption is replaced by the conjunction
  (REP)~+~(DISP)~+~(COAL) plus the embedding compatibility
  assumption plus the HHI threshold itself.
\end{itemize}

The universal model-free Concentration-Gap conjecture
of~\citep{lasser2026micro} remains open: a model-free
correlation-of-pressure-with-exploitation result is not provable
from current GFM machinery without a formal optimizer/selection
model, and the present paper's scoped discharge requires three
explicit structural conditions. The reverse direction of
Theorem~\ref{thm:concgap-scoped} also operates only at the
proxy-truth-norm level, not the alignment-property level: a
$g$-level lower bound would require additional inverse-Lipschitz
structure on $g$ that the GFM apparatus does not provide.
\S\ref{sec:discussion} states what each open item would require.

\paragraph{Organization.} \S\ref{sec:setup} fixes shared notation,
inheriting from the GFM
sequence~\citep{lasser2026micro,lasser2026paper8a,lasser2026paper8b,lasser2026exo,lasser2026paper10}.
\S\ref{sec:bounded_coev} proves
Theorem~\ref{thm:c7-corollary}. \S\ref{sec:concgap} proves
Theorem~\ref{thm:concgap-scoped}. \S\ref{sec:composition} locates
where each theorem enters
\citep[Deployment Safety]{lasser2026paper10}'s deployment-claim
hypothesis set. \S\ref{sec:audits} specifies the operational
audit procedures for (C7.RATE), (REP), (DISP), (COAL), and the
embedding compatibility precondition.
\S\ref{sec:discussion} collects the open items.
